% 替换位置：第6.5节起至第6章结束；整体替换原6.5--6.8
% 这是修改片段，不是完整文档；不要整份直接追加到原文末尾。
\subsection{多起点构造与模拟退火邻域搜索}
\subsubsection{可行方案构造与资源解码}
首先利用多起点随机扰动构造完整可行方案，再将每个架次的货箱集合、机型及批次顺序作为搜索状态。
随机性仅用于搜索，不表示需求、天气或设备状态具有随机变化。
对每个候选架次枚举所涉及服务区的访问顺序，并重新检查质量、体积、逐段载荷、能耗和硬时限。
同型号机体与电池分别维护可用时间，架次最早准备时刻为
\begin{equation}
s_k^{\mathrm{earliest}}=\max\{r_{u_k}^{\mathrm{body}},r_{e_k}^{\mathrm{battery}}\}.
\end{equation}
机体返航后恢复可用，电池返航后充电，只有充满才允许下一次准备。
候选方案改变后重新解码后续资源与任务时刻，不沿用已经失效的资源安排。

\subsubsection{邻域与候选接受}
邻域操作包括机型变更、批次插入或交换、货箱移动或交换、批次合并和拆分，
以及两个批次合并后重新分配货箱或按目的地重组。
以完成时间和延误为主的辅助评价函数驱动模拟退火，允许暂时接受部分较差候选以离开局部解；
最终保留方案仍按式~\eqref{eq:q2_objective}的字典序目标筛选。
不满足硬约束的候选不进入接受比较，不能用惩罚项放宽医疗和首批截止时间。

\subsubsection{搜索设置与范围}
邻域搜索采用20260925、20260926和20260927三个随机种子，每组搜索时间预算为$180\,\mathrm{s}$。
每架次访问点数上限取3或4，通过不同搜索设置构造候选方案。
按照字典序目标筛选满足第二问约束的候选，并对选定方案独立核验。
访问点数上限、邻域和壁钟时间均为计算限制，不是原题约束；未穷尽全部可行调度。

\subsection{调度结果与分析}
\subsubsection{基准方案汇总}
所得方案完成80箱按期交付。16箱医疗物资与30箱首批保障物资均满足各自硬时限，
两类标记允许重叠。结果见表~\ref{tab:q2_summary}。
\begin{table}[htbp]
\centering
\small
\caption{第二问基准调度结果}
\label{tab:q2_summary}
\begin{tabular}{lr}
\toprule
指标 & 结果 \\
\midrule
交付箱数 & 80 \\
加权延误$D$ & 0 \\
运输架次数 & 24 \\
多服务区架次数 & 6 \\
最晚返航时刻/s & 6208.684865 \\
总能耗/kWh & 69.240573590 \\
实际运输机/电池数量 & 8/14 \\
最低返航SOC & $20.958111\%$ \\
最小硬时限余量/s & 194.866445 \\
\bottomrule
\end{tabular}
\end{table}
A、B、C型分别执行12、6、6架次。各实体机均执行3架次，U08的末次返航决定$C_{\max}$。
该方案的最晚返航时间为$103.478081\,\mathrm{min}$，全部运输任务在给定库存内完成。
第一问的资源与时限约束、路线限制和目标顺序均与本问不同，不能仅比较18与24架次判断优劣。
\begin{table}[htbp]
\centering
\small
\caption{所得方案的实体无人机任务分配}
\label{tab:q2_units}
\begin{tabular}{llrrr}
\toprule
无人机 & 机型 & 架次数 & 累计占用/s & 末次返航/s \\
\midrule
U01 & A & 3 & 5390.093 & 5797.832 \\
U02 & A & 3 & 5699.660 & 5817.252 \\
U03 & A & 3 & 5215.296 & 5215.296 \\
U04 & A & 3 & 5337.313 & 6113.782 \\
U05 & B & 3 & 6024.529 & 6024.529 \\
U06 & B & 3 & 5418.168 & 5418.168 \\
U07 & C & 3 & 6195.854 & 6195.854 \\
U08 & C & 3 & 6208.685 & 6208.685 \\
\bottomrule
\end{tabular}
\end{table}

\subsubsection{多点架次的载荷与时间示例}
选取C型无人机U08执行的P009架次，其使用电池C-BAT-03，携带8箱物资，路线为
\begin{equation}
\mathrm{O01}\rightarrow\mathrm{S006}\rightarrow\mathrm{S007}\rightarrow\mathrm{O01}.
\end{equation}
初始载荷$74\,\mathrm{kg}$，在S006交付4箱后降至$37\,\mathrm{kg}$，
在S007交付其余4箱后空载返回。
\begin{table}[htbp]
\centering
\small
\caption{P009架次的逐航段载荷与能耗}
\label{tab:q2_multistop}
\begin{tabular}{lrrrr}
\toprule
航段 & 载荷/kg & 起始/s & 到达/s & 能耗/kWh \\
\midrule
O01$\rightarrow$S006 & 74 & 1968.159 & 2301.527 & 2.011156 \\
S006$\rightarrow$S007 & 37 & 2625.527 & 3023.843 & 1.040271 \\
S007$\rightarrow$O01 & 0 & 3347.843 & 3907.628 & 1.438274 \\
\bottomrule
\end{tabular}
\end{table}

该架次准备开始与起飞时刻分别为$1428.158932\,\mathrm{s}$和$1968.158932\,\mathrm{s}$，
两站交付完成时刻为$2625.527260\,\mathrm{s}$和$3347.843322\,\mathrm{s}$，
均在到达之后计入完整交接时间。
全架次能耗$4.489701013\,\mathrm{kWh}$，返航SOC为$43.878737\%$。
机体于$3907.628219\,\mathrm{s}$返航，但电池到$5956.922243\,\mathrm{s}$才充满，
表明机体与能源恢复可用时间必须分别处理。

\subsection{目标优先级对比与最优性边界}
\subsubsection{目标优先级比较}
在本文搜索候选池内，分别按$D\rightarrow C_{\max}\rightarrow E\rightarrow N$
和$C_{\max}\rightarrow D\rightarrow E\rightarrow N$筛选。
以$10^{-7}\,\mathrm{s}$作为完成时间的数值比较容差，两者选中相同方案。
\begin{table}[htbp]
\centering
\small
\caption{第二问候选池的目标优先级比较}
\label{tab:q2_priority}
\begin{tabular}{lrrrr}
\toprule
排序 & 加权延误 & 完成时间/s & 能耗/kWh & 架次 \\
\midrule
$D\to C_{\max}\to E\to N$ & 0 & 6208.684865 & 69.240573590 & 24 \\
$C_{\max}\to D\to E\to N$ & 0 & 6208.684865 & 69.240573590 & 24 \\
\bottomrule
\end{tabular}
\end{table}

两种排序得到相同结果仅反映所搜索范围，不证明及时性与完成时间在全部可行方案中没有冲突。

\subsubsection{最优性边界}
由于$D\geq0$且当前方案达到$D=0$，第一层目标达到其下界；
后续完成时间、能耗与架次数没有全局最优证明。
因此本问应表述为独立验证通过的当前搜索最好可行方案，不能仅凭多轮随机搜索声称全局最优。

\subsection{独立事件回放与可行性检查}
独立程序重新读取原始货箱、设备和时限，采用40位十进制精度复算能耗与时刻，
并按准备开始、返航充电、充满恢复可用三类事件进行回放，共72个运输资源事件。
检查80箱无遗漏、无重复、无错投，质量、体积、逐段载荷、SOC及全部硬时限均合规，
机体占用和电池使用、充电、复用无冲突。
最低返航SOC为$20.958111\%$，最小硬时限余量$194.866445\,\mathrm{s}$，
最大架次能耗重算差约$4.595\times10^{-16}\,\mathrm{kWh}$。
时间、能耗、质量、体积与SOC容差分别为$10^{-7}\,\mathrm{s}$、$10^{-8}\,\mathrm{kWh}$、
$10^{-9}\,\mathrm{kg}$、$10^{-12}\,\mathrm{m^3}$及$10^{-10}$。
判定使用未舍入数据，同刻先释放后占用。
全部24架次、80箱交付和电池充电记录另附；第二问的运输验证不能替代第三问的通信检查。
